\ssr{ВВЕДЕНИЕ}

\textbf{Цель работы:}
разработать и реализовать программное обеспечение на языке Python для извлечения и валидации данных из текстовых файлов, полученных из PDF-документов с применением библиотеки \texttt{PyPDF2} и регулярных выражений. Достижение цели требуется выполнить с использованием генеративных возможностей не менее трёх больших языковых моделей (БЯМ) для сравнительного анализа их эффективности в задачах кодинга.

\textbf{Поставленные задачи.}
\begin{enumerate}
    \item Выполнить аналитический обзор выбранных больших языковых моделей (БЯМ), включая локально разворачиваемые решения.
    \item Спроектировать и разработать регулярные выражения для решения задачи по варианту (валидация пробельных символов).
    \item Реализовать функцию поиска подстрок в PDF-файле, возвращающую кортеж вида:
    \begin{itemize}
        \item булево значение (истина/ложь), определяющее наличие искомого паттерна;
        \item список кортежей с найденными вхождениями и их координатами (номер страницы и номер строки документа).
    \end{itemize}
    \item Разработать программное обеспечение, принимающее на вход путь к PDF-файлу и использующее реализованную функцию.
    \item Задокументировать процесс взаимодействия с каждой БЯМ: привести промпты, полученные результаты, итерации исправлений ошибок и рефакторинга.
    \item Провести проверочные запуски на тестовом наборе файлов и сформировать таблицу с результатами (признак нахождения, координаты).
    \item Выполнить исследовательский анализ стабильности генерации кода и влияния режимов рассуждений (Reasoning) на качество алгоритмов.
\end{enumerate}

\textbf{Вариант:}
\begin{itemize}
    \item \textbf{задача} --- проверка отсутствия неразрывного пробела (поиск ошибок вёрстки);
    \item \textbf{объект поиска} --- специальные сочетания слов и чисел, где обычный пробел недопустим во избежание разрыва строки;
    \item \textbf{ключевые паттерны} --- комбинации вида: <<рисунок 1>>, <<таблица 2>>, числовые диапазоны (<<от 1 до 5>>), сравнения (<<менее 5>>, <<свыше 25>>);
    \item \textbf{инструментарий} --- язык Python, библиотека \texttt{PyPDF2}, модуль \texttt{re}.
\end{itemize}

\clearpage